\documentclass{article}
\usepackage{amsmath,amssymb}
\usepackage{xurl}
\usepackage[hidelinks]{hyperref}
% Shared commands for notes in docs/paper-gaps/.

\DeclareUnicodeCharacter{2080}{\ifmmode{}_{0}\else\textsubscript{0}\fi}
\DeclareUnicodeCharacter{2081}{\ifmmode{}_{1}\else\textsubscript{1}\fi}
\DeclareUnicodeCharacter{2082}{\ifmmode{}_{2}\else\textsubscript{2}\fi}
\DeclareUnicodeCharacter{2099}{\ifmmode{}_{n}\else\textsubscript{n}\fi}

\newcommand{\Lean}{\textsc{Lean}}
\ExplSyntaxOn
\NewDocumentCommand{\leanid}{m}{
  \ifmmode
    \textsf{\detokenize{#1}}
  \else
    \group_begin:
    \urlstyle{sf}
    \tl_set:Nn \l_tmpa_tl {#1}
    \tl_replace_all:Nnn \l_tmpa_tl {\_} {_}
    \exp_args:NV \path \l_tmpa_tl
    \group_end:
  \fi
}
\ExplSyntaxOff
\providecommand{\tr}{\operatorname{tr}}
\newcommand{\tnleanIssueBase}{https://github.com/LionSR/TNLean/issues/}
\newcommand{\tnleanPrBase}{https://github.com/LionSR/TNLean/pull/}
\newcommand{\tnleanDocBase}{https://sirui-lu.com/TNLean/docs/}
\newcommand{\ghissue}[1]{\href{\tnleanIssueBase#1}{GitHub issue \##1}}
\newcommand{\ghpr}[1]{\href{\tnleanPrBase#1}{PR \##1}}
\newcommand{\doclink}[1]{\href{\tnleanDocBase#1.html}{\path{#1}}}

% Dirac notation and the identity operator, as in blueprint/src/macros/common.tex.
% \providecommand keeps note-local definitions authoritative.
\usepackage{dsfont}
\providecommand{\ket}[1]{|#1\rangle}
\providecommand{\bra}[1]{\langle #1|}
\providecommand{\braket}[2]{\langle #1 | #2 \rangle}
\providecommand{\ketbra}[2]{|#1\rangle\!\langle #2|}
\providecommand{\Id}{\mathds{1}}
\providecommand{\one}{\mathds{1}}
\providecommand{\C}{\mathbb{C}}
\providecommand{\MN}[1]{M_{#1}(\C)}
\providecommand{\MD}{M_D(\C)}

% Verdict marker: \gapnote{<kind>}{<status>}. Kind is the policy
% classification (clarification, local-correction, scope-restriction,
% unfaithful, false-source, open-gap); status is open, wip (elimination
% underway), resolved, or historical. Machine-read by the site index;
% prints a verdict line.
\newcommand{\gapnote}[2]{\par\noindent\textbf{Verdict:} #1 (#2).\par}

\title{Finite dual-path deformation of regular PEPS flux strings}
\author{}
\date{2026-10-05}
\begin{document}
\maketitle
\gapnote{scope-restriction}{open}

\section{Source and orientation}
The source is Schuch, Cirac, and P\'erez-Garc\'ia,
\href{https://arxiv.org/abs/1001.3807}{arXiv:1001.3807},
Definition 6.13 and Lemma 6.14, local source
\path{Papers/1001.3807/paper_v3.tex}, lines 2174--2214.
The endpoint measurement is Theorem 6.15, lines 2216--2267;
lines 2230--2250 explain the four-site measurement and its independence
from the outgoing virtual indices. The figures
\path{figs5/def-flux-pair.pdf} and
\path{figs5/detect-flux-end.pdf} fix the native orientation.
This note belongs to \ghissue{8703}, under \ghissue{8271}.

Horizontal native arrows point right and vertical arrows point down.
The source's eastward dual segment puts $U_g$ on its crossed downward
bond; a southward dual segment puts $U_{g^{-1}}$ on its crossed
rightward bond. The conjugacy labels of the initial and final plaquettes
are consequently $[g]$ and $[g^{-1}]$ in the clockwise convention.
The pre-existing graph detector follows the opposite, counterclockwise
plaquette walk. It reads the inverse classes. This is an orientation
change, not an equality between a conjugacy class and its inverse.

\section{Paths and the derived gauge}
Let $X=\mathbb Z_w\times\mathbb Z_h$ with $w,h>0$.
A plaquette is labelled by its lower-left primal vertex.
A finite dual path is an actual sequence of east, west, north and south
steps in the standard path construction of a directed graph. Distinct
native bond labels remain distinct at periods one and two. Let
$n(p)=(n_h(p),n_v(p))$ be the signed integer crossing counts.
An east step from $(x,y)$ contributes $(0,\delta_{(x+1,y)})$;
a north step contributes $(\delta_{(x,y+1)},0)$; reversed steps
contribute their negatives. Crossings add with multiplicity.

For $R\subseteq X$, the relation $p\sim_R q$ is generated by adjacent
backtrack removal and the exchange of the east--north and north--east
sides of a dual square enclosing a vertex of $R$. It is closed under
symmetry, transitivity, reversal and concatenation. It is defined
without any premise about a PEPS state, a connection, a gauge, or flatness.
Induction on these moves derives an integer potential $f:X\to\mathbb Z$
vanishing outside $R$ with
\begin{align}
    n_h(q)(v)-n_h(p)(v) & =f(v+(1,0))-f(v), \\
    n_v(q)(v)-n_v(p)(v) & =f(v)-f(v+(0,1)).
    \label{eq:flux_gap_gradient}
\end{align}
For the positive square move, $f$ is the point mass at the enclosed
primal vertex. Reversing a move negates its potential and composing
moves adds potentials. Retracing a whole path and removing a dual-square
loop are consequences of these generators.

For any group $G$ and $g\in G$, the inserted native group labels are
$u_h^p(v)=g^{n_h(p)(v)}$ and $u_v^p(v)=g^{n_v(p)(v)}$.
Taking $k(v)=g^{f(v)}$ in (\ref{eq:flux_gap_gradient}) derives
\begin{align}
    u_h^q(v) & =k(v+(1,0))u_h^p(v)k(v)^{-1}, \\
    u_v^q(v) & =k(v)u_v^p(v)k(v+(0,1))^{-1}.
    \label{eq:flux_gap_gauge}
\end{align}
All products in this calculation are powers of the same element;
$G$ need not be abelian.

\section{Actual physical equality and the exterior}
For a finite-dimensional matrix representation $U$ and site tensors
$A_v$ which are virtually $G$-invariant on $R$, the gauge in
(\ref{eq:flux_gap_gauge}) is absorbed in the actual native bond-end sum.
The physical coefficient is unchanged for every physical configuration.
The site tensors may be inhomogeneous, their physical alphabets may vary,
and site tensors outside $R$ are unrestricted. No regularity, unitarity
or injectivity is needed for this deformation argument beyond virtual
invariance at the swept sites.

The exterior statement has an exact boundary. On every bond incident
to $R$, use the path's prescribed matrix $U_{u_e^p}$ or $U_{u_e^q}$.
On every bond whose two endpoints lie outside $R$, retain an arbitrary
fixed matrix $E_e$. These exterior matrices need not be representation
matrices and need not commute with $U_g$. Since $k=1$ at both ends of
such a bond, its matrix is unchanged. Crossing bonds are part of the
vertex star and receive the derived gauge change. This proves equality
of the actual closed contractions, rather than assuming a state equality
or asserting equality for fixed open boundary columns.

The main declaration is
\leanid{TNLean.PEPS.TorusDualHomotopy.torusBondNetwork_eq_with_exterior}.
Its geometric input is
\leanid{TNLean.PEPS.TorusDualHomotopy.exists_potential};
the contraction argument is
\leanid{TNLean.PEPS.torusBondNetwork_groupGauge_of_invariant_on}.
The state-valued specialization is
\leanid{TNLean.PEPS.TorusDualHomotopy.torusDualFluxState_eq}.

\section{Endpoint flux and measurement conventions}
For native labels $u=(u_h,u_v)$, clockwise holonomy is
\begin{align}
    H_u(x,y) & =u_h(x,y)^{-1}u_v(x+1,y)u_h(x,y+1)u_v(x,y)^{-1}.
\end{align}
The signed curl is the sum of the top and right crossing counts minus
the bottom and left crossing counts. The single-step calculation
telescopes along every path $p:a\to b$:
\begin{align}
    (\operatorname{curl}n(p))(z) & =\delta_a(z)-\delta_b(z),     \\
    H_{u^p}(z)                   & =g^{\delta_a(z)-\delta_b(z)}.
    \label{eq:flux_gap_endpoints}
\end{align}
Thus distinct endpoints carry $g,g^{-1}$ and all other plaquettes
have identity holonomy. If $a=b$, the holonomy is identity everywhere,
even for a winding loop. Local absence of flux does not identify its
global state with the state having no inserted string.

For $w,h\geq3$, the native regular-matrix contraction is identified
entrywise with the ordered-edge regular graph contraction, including
at both seams. The fixed complete projective measurement on the four
original physical spins at $z$ therefore reads $[H_{u^p}(z)^{-1}]$
in the counterclockwise convention. Inversion gives the canonical
bijection $\iota([s])=[s^{-1}]$ of conjugacy classes. Relabelling the
outcomes by $Q_C^{\rm cw}=Q_{\iota(C)}^{\rm ccw}$ gives the source's
$[g]$ and $[g^{-1}]$ at the initial and final plaquettes. The extra
orthogonal outcome is unchanged. All outcomes remain positive,
orthogonal and complete; the projectors are independent of $p$ and $g$.

The source-oriented measurement is
\leanid{TNLean.PEPS.IsGIsometric.exists_torusDualFlux_clockwise_cutMeasurement}.
It acts on the cut matrix formed from the actual inserted contraction.
Regular G-isometry also implies that every such native string state and
every physical cut matrix are nonzero, by the nonvanishing theorem for
regular twisted graph states and the entrywise contraction identity.
The input to endpoint detection is therefore proved nonzero. In the cyclic group $C_3$,
a generator and its inverse lie in different conjugacy classes. That
example distinguishes the two conventions; $S_3$ alone cannot, since
every permutation in $S_3$ is conjugate to its inverse.

\section{Local parent consequences}
Suppose $w,h\geq3$ and the original homogeneous tensor is virtually
$G$-invariant. Identity holonomy on one plaquette gives a gauge supported
on its four sites which removes the four internal insertions. Every
complementary physical slice of the actual inserted state then lies in
the original plaquette PEPS range. Hence every positive parent interaction
with that range as kernel annihilates the state there, including the
canonical orthogonal projector onto the range's orthogonal complement.
Substituting (\ref{eq:flux_gap_endpoints}) derives these constraints at
every nonendpoint plaquette, with no assumed vanishing curl or holonomy
in the final path statement. A closed loop satisfies all such local
parent constraints. This is a range and local-annihilation statement;
it does not assert equality of all reduced density matrices away from
the endpoints.

\section{Remaining scope and continuation}
The established equality is for finite paths related by the specified
local moves on a closed labelled torus. The lifted integer displacement
reduces modulo the two periods to the endpoint difference and is
invariant under those moves. Thus paths with the same periodic
endpoints but different winding are not identified. The result does not
classify all generated homotopy classes, prove that every equal-winding
pair is related, or turn endpoint agreement alone into equality of states.
The source lemma's deformation language must be read with its topology
and allowed local deformations specified.

The geometry and contraction identity hold for every positive period.
The four-site detector and local parent consequences use a simple graph
and require both periods at least three. A labelled-graph treatment of
loops and parallel bonds is needed to extend those physical statements
to smaller periods. Other surfaces and open boundary conditions require
a corresponding path geometry and explicit boundary contraction theorem.
No unrestricted version of Lemma 6.14 is marked as proved by this packet.
The continuation under \ghissue{8703} and \ghissue{8271} is to extend these
precisely stated boundaries when the wider source setting is required,
rather than introduce a flatness assumption or a supplied gauge witness
into the geometric deformation theorem.

\section{Finite rectangular continuation}
A finite rectangular specialization has passed strict Lean compilation and
axiom checks in draft \ghpr{8726}.
Choose a translated lifted rectangle $[0,W]\times[0,H]$ with $W<w$ and
$H<h$, retain the four unit-step labels, and require both paths to have
the same lifted endpoints. A bottom-row and column comb reduces the
paths by adjacent-square exchanges and backtrack cancellations.
The permitted swept primal sites form the whole rectangular interior
$R=\{(x_0+i+1,y_0+j+1):i<W,\ j<H\}$, with coordinates reduced modulo
the two periods. The contraction statement requires virtual invariance
at every site of this chosen $R$. Arbitrary fixed exterior matrices are
allowed only on bonds whose two endpoints both lie outside $R$.

This specialization does not assert deformation avoiding an arbitrary
region, equality of winding classes, an open-boundary theorem, or equality
with a vacuum reduced density matrix. The unrestricted source gap remains open.

\section{Collared open-boundary continuation}
Draft \ghpr{8729} develops a further specialization with a one-site collar:
$W+1<w$ and $H+1<h$. Its proposed open coefficient is an independent finite
sum over region incidences, with boundary indices fixed on crossing bonds
and matrices on internal labelled bonds. A delta-exterior completion is
intended to identify this sum with a closed contraction exactly.
The collar makes a derived supported gauge trivial at crossing endpoints,
which would give coefficient equality and contraction with an arbitrary
joint boundary tensor. This continuation is under strict verification;
these statements are not yet reported as completed here. The unrestricted
source gap, arbitrary-region avoidance, cross-winding equality and vacuum
reduced-density conclusions remain open.

\end{document}
